% =====================================================================
%  paper.tex  --  arXiv preprint (journal-flexible)
%
%  "A Verified Floating-Point Dijkstra in Lean 4 + Rust"
%
%  Compile with:
%      latexmk -pdf paper.tex
%  or:
%      pdflatex paper && bibtex paper && pdflatex paper && pdflatex paper
% =====================================================================
\documentclass[11pt,a4paper]{article}

% --- packages --------------------------------------------------------
\usepackage[T1]{fontenc}
\usepackage[utf8]{inputenc}
\usepackage{lmodern}
\usepackage[english]{babel}
\usepackage[margin=1.1in]{geometry}
\IfFileExists{microtype.sty}{\usepackage{microtype}}{}
\usepackage{amsmath,amssymb,amsthm}
\IfFileExists{mathtools.sty}{\usepackage{mathtools}}{}
% Conditional package loading: load only when the package and all of
% its hard dependencies are present in this TeX install.
\IfFileExists{xcolor.sty}{\usepackage{xcolor}}{}
\IfFileExists{hyperref.sty}{%
  \IfFileExists{etoolbox.sty}{%
    \usepackage[hidelinks]{hyperref}%
  }{}%
}{}
\IfFileExists{cleveref.sty}{\usepackage{cleveref}}{}
\IfFileExists{enumitem.sty}{\usepackage{enumitem}}{}
\IfFileExists{listings.sty}{\usepackage{listings}}{}
\IfFileExists{booktabs.sty}{\usepackage{booktabs}}{}
\usepackage{xspace}
\usepackage{url}

% Fallbacks for commands provided by the optional packages above.
% \providecommand is a no-op if the command was already defined.
\providecommand{\textcolor}[2]{#2}
\providecommand{\color}[1]{}
\providecommand{\href}[2]{#2}
\providecommand{\cref}[1]{Section~\ref{#1}}
\providecommand{\Cref}[1]{Section~\ref{#1}}

\makeatletter
\@ifpackageloaded{enumitem}{}{%
  % Strip optional argument from enumerate/itemize when enumitem is missing.
  \let\enitfb@enumerate\enumerate
  \let\enitfb@endenumerate\endenumerate
  \let\enitfb@itemize\itemize
  \let\enitfb@enditemize\enditemize
  \renewenvironment{enumerate}[1][]{\enitfb@enumerate}{\enitfb@endenumerate}%
  \renewenvironment{itemize}[1][]{\enitfb@itemize}{\enitfb@enditemize}%
}
\makeatother

% --- listings: Rust + Lean (only configured if package available;
%               currently unused in the drafted text but kept for
%               future code excerpts in section stubs) ---------------
\makeatletter
\@ifpackageloaded{listings}{%
  \lstdefinelanguage{Rust}{%
    morekeywords={fn,let,mut,pub,struct,enum,trait,impl,for,while,loop,if,else,
                  match,return,use,mod,crate,as,in,where,unsafe,move,ref,Self,self,
                  true,false,Some,None,Ok,Err,break,continue,async,await,dyn,extern,
                  type,const,static},%
    morekeywords=[2]{u8,u16,u32,u64,usize,i8,i16,i32,i64,isize,f32,f64,bool,char,str,String,Vec,Option,Result,Box},%
    morecomment=[l]{//},%
    morecomment=[s]{/*}{*/},%
    morestring=[b]"%
  }
  \lstdefinelanguage{Lean}{%
    morekeywords={def,theorem,lemma,example,structure,inductive,class,instance,
                  where,match,with,fun,let,if,then,else,do,return,by,abbrev,
                  section,end,namespace,open,variable,variables},%
    morekeywords=[2]{Prop,Type,Sort,True,False,And,Or,Not,Iff,Exists,Forall,Nat,Int,Float,Bool,List,Option,Fin,WithTop,NNReal,Float},%
    morecomment=[l]{--},%
    morecomment=[s]{/-}{-/},%
    morestring=[b]"%
  }
  \lstset{basicstyle=\ttfamily\small,numbers=none,breaklines=true,frame=single}%
}{}
\makeatother

% --- macros ----------------------------------------------------------
\newcommand{\Lean}{Lean~4\xspace}
\newcommand{\Mathlib}{\textsc{Mathlib}\xspace}
\newcommand{\BMSSP}{\textsc{BMSSP}\xspace}
\newcommand{\Dijkstra}{Dijkstra\xspace}
\newcommand{\Bigo}{\mathcal{O}}
\newcommand{\dhat}{\hat{d}}
\newcommand{\dist}{\mathit{dist}}
\newcommand{\code}[1]{\texttt{#1}}
\newcommand{\file}[1]{\texttt{#1}}
% Allow line breaks inside long typewriter paths.
\sloppy
\hyphenpenalty=2500\tolerance=4000\emergencystretch=3em
% \todo macro: simple bold text in brackets. (When xcolor is available
% the manuscript build may grep for [TODO: and recolour by hand.)
\newcommand{\todo}[1]{\textbf{[TODO: #1]}}

\theoremstyle{plain}
\newtheorem{theorem}{Theorem}[section]
\newtheorem{lemma}[theorem]{Lemma}
\newtheorem{proposition}[theorem]{Proposition}
\newtheorem{corollary}[theorem]{Corollary}
\theoremstyle{definition}
\newtheorem{definition}[theorem]{Definition}
\newtheorem{remark}[theorem]{Remark}
\newtheorem{example}[theorem]{Example}

% --- title -----------------------------------------------------------
\title{\textbf{A Verified Floating-Point \Dijkstra\\
       in \Lean{} and \textsc{Rust}}\\
       \large{}with a Bridge-Axiom Inventory toward Zero-Trust Refinement}

\author{
  [Author 1]\and [Author 2]%
  \thanks{Repository: \texttt{https://github.com/ilerik/graphs-bench}}
}

\date{June 15, 2025}

% =====================================================================
\begin{document}
\maketitle

\begin{abstract}
\noindent
\Dijkstra's single-source shortest-path algorithm has been mechanically
verified in proof assistants for over twenty years, in eight or more systems
ranging from \textsc{Mizar} and \textsc{ACL2} to \textsc{Isabelle/HOL},
\textsc{Coq/VST}, and \textsc{Why3}. Yet \emph{every} published mechanization
either treats edge weights as mathematical reals or unbounded integers
(idealised models) or treats them as bounded machine integers with an explicit
overflow precondition, as in the most recent work of Mohan, Leow, and Hobor
(\emph{CAV 2021}). To our knowledge, no mechanized correctness proof of
\Dijkstra{} addresses \emph{IEEE-754 floating-point} edge weights, which is
the model every shipped Rust, C, and C++ implementation actually uses.

We close this gap. We present a verified, executable \Dijkstra{} in
\Lean{}~4 plus \Mathlib{}, refined to a real Rust implementation with
compressed-sparse-row (CSR) adjacency and \code{f64} edge weights. The Lean
side is organised as a three-layer architecture: a non-computable
\emph{specification} layer over \Mathlib's \code{SimpleGraph} hierarchy and
the metric \code{trueDist} valued in \code{WithTop NNReal}; a computable
\emph{algorithm} layer~\code{Sssp.Algo.Dijkstra} carrying the principal
correctness theorem~\code{dijkstra\_correct}; and an
\emph{operational refinement} layer~\code{Sssp.Refine.Dijkstra} that mirrors
the Rust source line-for-line. The bridge between layers is mediated by two
named axiom blocks --- \code{NumericBridge} (21 IEEE-754 facts) and
\code{HeapBridge} (one lazy-heap simulation axiom, with three conditional
discharge routes already proved) --- isolating exactly what a
\Lean{}~$\leftrightarrow$~IEEE-754~$\leftrightarrow$~\textsc{Rust} pipeline
must currently assume; \code{RelaxBridge} and \code{GraphBridge} were
originally axiomatic and have since been fully proved, in the absence of a
port of \textsc{Flocq} to \Lean{}~4. The principal theorem
\code{refine\_dijkstra\_correct} establishes that, modulo these two named
axiom blocks,
the Rust \Dijkstra{} computes \code{trueDistNat} on every
\code{ValidRustGraph}; an executable fixture harness shared between Lean
\code{native\_decide} and \code{cargo} JSON tests provides regression
coverage. We give a roadmap (Phase~3c) for discharging each axiom, and
an empirical study comparing \Dijkstra{} to the recently introduced
$\Bigo(m\cdot \log^{2/3} n)$ algorithm of Duan, Mao, Mao, Shu and
Yin~\cite{DuanMaoMaoShuYin2025} that motivates why a verified textbook
\Dijkstra{} is the practical baseline.
\end{abstract}

% =====================================================================
\section{Introduction}\label{sec:introduction}

\paragraph*{A 65-year-old algorithm, still without a floating-point proof.}
\Dijkstra's single-source shortest-path algorithm~\cite{Dijkstra1959} is, by
any measure, one of the most-mechanically-verified algorithms in
computer science. Over the past twenty years it has been formalised in
\textsc{Mizar}~\cite{Chen2003}, \textsc{ACL2}~\cite{MooreZhang2005},
\textsc{Jahob}~\cite{MangeKuhn2007}, \textsc{KeY}~\cite{Klasen2010},
\textsc{Why3}~\cite{Filliatre2011Toccata}, an SMT-based bounded
checker~\cite{LiuNagelTaghdiri2012}, \textsc{Isabelle/HOL} via the Refinement
Framework~\cite{NordhoffLammich2012,LammichNipkow2019},
\textsc{Dafny}~\cite{Leino2023ProgramProofs}, and most recently, in
\textsc{Coq} on top of \textsc{CompCert}~\cite{Leroy2009CompCert} and
\textsc{VST}~\cite{AppelVST2011} as part of the \textsc{CertiGraph}
project~\cite{WangCaoMohanHobor2019,MohanLeowHobor2021}. A reasonable
practitioner could be forgiven for thinking the topic exhausted.

It is not. Every one of the developments above either
\emph{(i)}~works over an idealised arithmetic model in which edge weights are
unbounded integers or mathematical reals (most of the proof-assistant
community: \textsc{Mizar}, \textsc{ACL2}, \textsc{Jahob}, \textsc{KeY},
\textsc{Why3}, \textsc{Isabelle/HOL}, \textsc{Dafny}), or
\emph{(ii)}~retains machine integers but takes care to specify an
overflow precondition (\textsc{CertiGraph}: edge weights bounded by
$\lfloor \mathrm{INT\_MAX}/\mathrm{size} \rfloor$, after Mohan et al.\ uncovered
that the more obvious bound $\lfloor \mathrm{INT\_MAX}/(\mathrm{size}-1)\rfloor$
is unsound~\cite[\S 3.2]{MohanLeowHobor2021}). Mohan et al.\ explicitly call
out the missing third option:
\begin{quote}
\itshape``Sedgewick's sidestep [to \code{double} weights] entails enduring
the inevitable round-off intrinsic to floating-point arithmetic, and so his
algorithm computes approximate optimal costs rather than exact ones.
Sedgewick does not specify any bounds on input edge weights, and accordingly
does not (and cannot) provide any bound on this accumulated
error.''~\cite[\S 3.3]{MohanLeowHobor2021}
\end{quote}
This is the exact gap our work occupies: a verified \Dijkstra{} over
IEEE-754 \code{f64} weights, with a real Rust implementation, refined to a
correctness theorem in \Lean.

\paragraph*{Why now?}
The gap matters for two further reasons beyond completeness of the
mechanization landscape.

First, after fifty years of stability, the SSSP problem has acquired a
genuinely new asymptotic frontier. Duan, Mao, Mao, Shu and
Yin~\cite{DuanMaoMaoShuYin2025} have given the first
$\Bigo(m \cdot \log^{2/3} n)$ deterministic algorithm for directed SSSP with
non-negative real weights in the comparison-addition model. This is the
first asymptotic improvement over \Dijkstra-with-Fibonacci-heap
$\Bigo(m + n \log n)$~\cite{FredmanTarjan1987} for sparse directed graphs.
Engineering claims about the new algorithm --- correctness, numerical
stability, where the asymptotic crossover with \Dijkstra{} actually lies ---
require a mechanized baseline against which to be judged. Without a verified
\Dijkstra{} over the same data structures and the same arithmetic, no such
comparison can be claimed at the level of foundational proof.

Second, Rust has emerged as the dominant systems language for
high-performance graph workloads. Yet despite the rapid maturation of Rust
verifiers in recent years (\textsc{Verus}~\cite{Lattuada2023Verus},
\textsc{Creusot}~\cite{DenisJourdanMarche2022Creusot},
\textsc{Aeneas}~\cite{HoProtzenko2022Aeneas},
\textsc{Prusti}~\cite{Astrauskas2019Prusti},
\textsc{Kani}~\cite{VanHattum2022Kani},
\textsc{RustHorn}~\cite{Matsushita2020RustHorn,Matsushita2022RustHornBelt}),
to our knowledge \emph{no verified shortest-path algorithm has been
published in any Rust verifier.} The Rust crate ecosystem (\code{petgraph},
\code{pathfinding}, \code{graphlib}) ships unverified \code{f64}-weighted
\Dijkstra{} as the default; a verified replacement is overdue.

\paragraph*{Why \Lean{} 4?}
\Lean{}~\cite{deMouraUllrich2021Lean4} together with \Mathlib~\cite{Mathlib2020}
gives the cleanest available formalisation of graph metrics:
\code{SimpleGraph.dist}, \code{SimpleGraph.edist}, walks, paths and
connectivity~\cite{MathlibSimpleGraph}. Crucially, however, \Mathlib{}
contains \emph{no} executable Dijkstra: the closest existing artefact is
\code{Quiver.shortestPath}~\cite{MathlibQuiver}, which is a non-computable
\code{Classical.choose}-style oracle on rooted connected quivers, and there
is no edge-weighted graph hierarchy with ordered semirings of weights at all.
A self-contained, executable, edge-weighted SSSP development in \Lean{} 4
fills a real combinatorial-algorithmic gap, complementing the deeply
developed but still mostly non-algorithmic side of \Mathlib's combinatorics
library.

\Lean{} 4 also gives us a clean Rust bridge:
\textsc{Aeneas}~\cite{HoProtzenko2022Aeneas,HoFromherzProtzenko2024Aeneas}
provides a sound functional translation of safe Rust into a pure
$\lambda$-calculus, with \Lean{} 4 among its supported back-ends. We do not
use \textsc{Aeneas}'s automatic translator in this work --- we instead
hand-translate the Rust source~\code{src/dijkstra.rs} into a mirror
\Lean{} model~\code{Sssp.Refine.Dijkstra} that we then prove against a
verified, abstract \Lean{} reference implementation~\code{Sssp.Algo.Dijkstra}
--- but the architecture is deliberately compatible with future automation
via the \textsc{Aeneas} pipeline.

\paragraph*{Contributions.}
This paper makes four concrete contributions, each grounded in the Lean
development at \code{formal/lean/Sssp/}:
\begin{enumerate}[leftmargin=2.2em,itemsep=0.35em]
\item \textbf{The first verified \Dijkstra{} in \Lean{} 4.}
      We give a total, computable function~\code{Sssp.Algo.dijkstra} together
      with a closed correctness theorem
      \[
        \code{dijkstra\_correct}\;:\;
        \forall G\,s\,v,\;
        \code{Algo.dijkstra}\,G\,s\,v
        = \code{trueDist}\,G\,s\,v
      \]
      where \code{trueDist} is the \Mathlib-style infimum over walks valued
      in $\mathrm{WithTop}\,\mathbb{R}_{\geq 0}$. The proof is split into
      soundness (every relaxation respects the metric) and completeness (a
      shortest walk has length at most $n$, so $n$ rounds of full edge
      relaxation suffice). See
      \code{formal/lean/Sssp/Algo/Dijkstra.lean:318}.

\item \textbf{The first \Lean{} refinement of an \code{f64}/CSR \Dijkstra{}
      against a real Rust implementation.}
      We expose a verified Rust binding via the theorem
      \begin{multline*}
        \code{refine\_dijkstra\_correct}\;:\;
        \forall\, (vg : \code{ValidRustGraph}\,n\,g)\,(s\,v : \mathrm{Fin}\,n),\\
        (\code{dijkstra}\,g\,s.\code{val})[v.\code{val}]!
          = \code{withTopNatToFloat}\,(\code{trueDistNat}\,vg.\code{toGraph}\,s\,v)
      \end{multline*}
      living at \code{formal/lean/Sssp/Refine/RefineCorrectness.lean:23},
      where \code{Sssp.Refine.Dijkstra} mirrors \code{src/dijkstra.rs}'s
      lazy-binary-heap implementation line-for-line and \code{trueDistNat}
      (\code{formal/lean/Sssp/Refine/NumericBridge.lean:40}) is an executable
      $\mathbb{N}$-shadow of \code{trueDist} obtained by a classical-choice
      nat preimage of the \code{NNReal} answer.

\item \textbf{A bridge-axiom inventory for the Lean--IEEE-754--Rust stack.}
      Where the proof currently relies on assumptions, we name and isolate
      them in two bridge modules:
      \begin{itemize}[leftmargin=1.2em]
      \item \code{NumericBridge}~--- 21 IEEE-754 facts about \code{Float}
            (commutativity, associativity on relevant subdomains,
            monotonicity, behaviour of \code{nnrealToFloat} on
            \code{ofNat}-style inputs);
      \item \code{HeapBridge}~--- the lazy-heap simulation axiom
            \code{dijkstraHeap\_eq\_dijkstraRelax}
            (\code{formal/lean/Sssp/Refine/HeapBridge.lean:23}), with three
            conditional discharge routes already proved:
            \code{dijkstraHeap\_eq\_dijkstraRelax\_of\_upper} (:27),
            \code{of\_complete} (:43), and \code{of\_edgeUpper} (:55), which
            reduce the standing axiom to one Dijkstra-completeness lemma.
      \end{itemize}
      \code{RelaxBridge} and \code{GraphBridge} are fully proved (no axioms).
      An axiom-elimination roadmap (Phase~3c) is laid out in
      \code{formal/FUTURE\_WORK.md}; this paper documents the current
      surface and explains why each axiom is the natural cut-point given
      \Mathlib's lack of a \textsc{Flocq}-style IEEE-754
      formalisation~\cite{BoldoMelquiond2011Flocq,BoldoMelquiond2017Book}.

\item \textbf{A regression harness shared between \Lean{} and Rust.}
      A common JSON fixture format --- four small directed graphs with
      pre-computed expected distances --- is checked by both
      \code{native\_decide}-driven \code{\#guard} statements in
      \code{Sssp.Fixtures.*} and \code{cargo test shared\_json\_fixtures}
      against \code{src/dijkstra.rs}, glued together by
      \code{formal/scripts/check-fixtures.sh}. The harness establishes
      operational agreement between the verified Lean model and the Rust
      binary on concrete inputs, complementing (but not replacing) the
      symbolic refinement proof.

\item \textbf{An empirical numerical-estimate study.}
      We benchmark \Dijkstra{} (binary heap, \code{f64} weights, CSR) against
      our Rust implementation of \BMSSP{} (the new
      $\Bigo(m\cdot \log^{2/3} n)$ algorithm of
      Duan~et~al.~\cite{DuanMaoMaoShuYin2025}) on Erdős--Rényi-style random
      directed graphs from $n=10^3$ to $n=4 \cdot 10^6$. The benchmark
      confirms the predicted $1/\log^{1/3} n$ ratio decay within each
      recursion-depth regime, but also shows that the practical crossover
      lies at $n \gtrsim 2^{216}$. This both motivates why a verified
      textbook \Dijkstra{} is the practical baseline of interest and
      bounds, empirically, the regime in which any future verification of
      \BMSSP{} would be worth deploying. See \cref{sec:benchmark}.
\end{enumerate}

\paragraph*{Non-goals.}
We are explicit about what this paper does \emph{not} attempt:
\begin{itemize}[leftmargin=1.6em,itemsep=0.25em]
\item We do not prove a forward-error bound on the \code{f64}
      computation in the style of
      Higham~\cite{Higham2002Accuracy}: the \code{NumericBridge} axiom
      block names the IEEE-754 facts we rely on, and an explicit
      $(n-1)\cdot\mathrm{ulp}$-style bound is deferred to a
      \textsc{Flocq}-port follow-up.
\item We do not verify the \BMSSP{} algorithm of
      Duan~et~al.~\cite{DuanMaoMaoShuYin2025} in \Lean. Our Rust
      implementation of \BMSSP{} is unverified; it appears here only as a
      benchmark baseline for the \code{f64} \Dijkstra{} we do verify, and as
      motivation for why a clean \Lean{}-side foundation matters going
      forward.
\item We do not verify a runtime cost bound on \Dijkstra{}. \Mathlib's
      \code{Asymptotics.IsBigO} machinery is in place, but is not connected
      to a runtime cost model in the style of Charguéraud and Pottier's
      time credits~\cite{CharguraudPottier2015UnionFind,GueneauChargueraudPottier2018ESOP};
      we leave a \Lean-side cost monad to future work.
\end{itemize}

\paragraph*{Outline.}
\Cref{sec:related} surveys related work in six clusters. \Cref{sec:prelim}
fixes notation and recalls just enough \Lean, \Mathlib{} and IEEE-754
background. \Cref{sec:architecture} describes the three-layer architecture
of the Lean development. \Cref{sec:algo} presents the abstract verified
\Dijkstra. \Cref{sec:refine} steps through the Refine layer and the
two remaining bridge axiom blocks (with \code{RelaxBridge} and
\code{GraphBridge} discharged since the initial development).
\Cref{sec:fixtures} describes the fixture harness.
\Cref{sec:benchmark} reports the empirical \Dijkstra/\BMSSP{} comparison.
\Cref{sec:discussion} discusses Phase~3c (zero-axiom \Dijkstra),
Aeneas-driven extraction, and the path to a verified Higham-style rounding
bound.

% =====================================================================
\section{Related work}\label{sec:related}

We organise the literature into six clusters. \S\ref{sec:rel-abs} covers the
long history of \Dijkstra-on-abstract-weights mechanizations.
\S\ref{sec:rel-real} addresses the considerably shorter history of
\Dijkstra-on-real-machine-models. \S\ref{sec:rel-graph-lean} situates our work
in the \Lean/\Mathlib{} graph-theory ecosystem. \S\ref{sec:rel-rust} covers
Rust verifiers. \S\ref{sec:rel-fp} surveys floating-point in proof assistants
and the (small) intersection with shortest-path correctness.
\S\ref{sec:rel-cost} closes with refinement, extraction and cost-monad work
that is methodologically close but, for this paper, deferred.

\subsection{Mechanized \Dijkstra{} over abstract weights}\label{sec:rel-abs}

The earliest mechanized correctness proof of \Dijkstra{} we are aware of is
Chen's 2003 \textsc{Mizar} development~\cite{Chen2003}, formalised over
abstract integer weights and following the textbook presentation. Two years
later, Moore and Zhang~\cite{MooreZhang2005} gave a self-contained
``proof pearl'' in \textsc{ACL2}, again on a finite weighted graph with
unbounded integer arithmetic. The Java verifier community contributed two
case studies: Mange and Kuhn~\cite{MangeKuhn2007} verified a Java
implementation in \textsc{Jahob}, and Klasen~\cite{Klasen2010} did so in
\textsc{KeY}. Liu, Nagel, and Taghdiri~\cite{LiuNagelTaghdiri2012} took a
bounded-SMT approach for the same class of programs, trading foundational
guarantees for automation. Filliâtre's \textsc{Why3} development for the
Toccata gallery~\cite{Filliatre2011Toccata} remains a particularly
pedagogical reference, with explicit invariants \code{inv\_src}, \code{inv},
\code{inv\_succ}; weights are non-negative \code{int}s.

The most algorithmically sophisticated efforts on the Isabelle side come
from the Refinement Framework community. Nordhoff and
Lammich~\cite{NordhoffLammich2012} verified \Dijkstra{} via stepwise data
refinement of a monadic, nondeterministic specification, ending in an
efficient implementation parametrised over the Isabelle Collections
Framework; the AFP entry has been continuously maintained since 2012.
Lammich and Nipkow~\cite{LammichNipkow2019} subsequently produced a
short, elegant proof using purely functional priority search trees, applied
to both \Dijkstra{} and Prim. Wimmer and Lammich~\cite{WimmerLammich2017}
formalised the sister APSP algorithm of Floyd--Warshall in the same
framework.

The common feature of every development in this paragraph is that
arithmetic is idealised: weights live in $\mathbb{N}$, $\mathbb{Z}$, or an
abstract ordered monoid; addition is exact; comparisons are decidable; no
overflow, no rounding. The textbook code can be transcribed essentially
verbatim, and the proof effort goes into the algorithmic invariants
(``popped vertices have globally optimal distances''). Leino's
\textsc{Dafny} treatment in \emph{Program Proofs}~\cite{Leino2023ProgramProofs}
fits the same pattern, with SMT automation taking the place of interactive
tactics. To our knowledge, none of these works have been ported to a
floating-point or even an explicitly bounded-integer setting.

\subsection{Mechanized \Dijkstra{} against a real machine model}\label{sec:rel-real}

The picture changes sharply when one moves to verified
\emph{implementations} (in C, Java, OCaml) running against a formal compiler
semantics. The \textsc{CertiGraph} project~\cite{WangCaoMohanHobor2019,Wang2019Thesis},
built on \textsc{CompCert}~\cite{Leroy2009CompCert} and the
Verified Software Toolchain~\cite{AppelVST2011}, has produced the most
extensive verified-graph-algorithm library in any proof assistant, with
modules for union-find, garbage collection, mark, copy, and the textbook
graph algorithms. Mohan, Leow, and Hobor's CAV~2021 paper~\cite{MohanLeowHobor2021}
gives the most detailed treatment to date: full functional correctness of C
implementations of \Dijkstra, Kruskal, and Prim, including a verified binary
heap with Floyd's bottom-up build and \code{decrease\_priority}. Crucially,
their work uncovers (and corrects) two subtle issues invisible in idealised
mechanizations:
\begin{enumerate}[leftmargin=1.6em,itemsep=0.25em]
\item Sedgewick-style edge-weight bounds of $\lfloor \mathrm{INT\_MAX}/(\mathrm{size}-1)\rfloor$ are
      \emph{insufficient} to prevent intermediate-value overflow on the
      ``probing edge'' $u \rightarrow v$ examined in the relaxation loop;
      the correct bound is $\lfloor \mathrm{INT\_MAX}/\mathrm{size}\rfloor$
      ~\cite[Fig.~2]{MohanLeowHobor2021}.
\item Textbook treatments of Prim's algorithm wrongly insist on a connected
      input; the standard pseudocode actually produces a minimum spanning
      forest on disconnected input without modification.
\end{enumerate}
Their related-work paragraph names FP weights as the open problem we
address~\cite[\S 3.3]{MohanLeowHobor2021}.

A methodological cousin of \textsc{CertiGraph} is Charguéraud's
\textsc{CFML}~\cite{Charguraud2011CFML}, which extracts \emph{characteristic
formulae} from OCaml programs and verifies them in \textsc{Coq};
\textsc{CFML} has been used for a range of textbook algorithms, including
some shortest-path examples in the manual, although we are not aware of a
published \Dijkstra{} verification at the level of detail of
Mohan~et~al.~\cite{MohanLeowHobor2021}.
\textsc{Iris}~\cite{JungRustBelt2018} and the broader separation-logic
community have produced a wide range of verified \emph{data-structure}
correctness proofs (binary heaps, skew heaps, Fibonacci heaps), which
together with a refinement to imperative
HOL~\cite{Lammich2019JAR,Lammich2019LLVM} or LLVM are the ingredients of an
end-to-end imperative verification.

\subsection{Graph theory and graph algorithms in \Lean{}/\Mathlib}\label{sec:rel-graph-lean}

\Mathlib's combinatorics hierarchy~\cite{Mathlib2020,MathlibSimpleGraph} is
deep on the structural side: \code{SimpleGraph} is the central abstraction,
with \code{Walk}, \code{Path}, connectivity, subgraph,
matching/coloring/Hamiltonian/planar/regularity (Szemerédi's regularity
lemma), girth, density, triangle-removal, and adjacency/incidence/Laplacian
matrix theory all in place. \code{SimpleGraph.dist} and
\code{SimpleGraph.edist} (the latter \code{WithTop}-valued) define the
graph metric as $\inf$ over walks, but are non-computable. The closest
existing artefact to an algorithmic shortest-path is
\code{Quiver.shortestPath}~\cite{MathlibQuiver}, which uses
\code{Classical.choose} to extract a path of minimal length on a rooted
connected quiver --- again non-computable, and with no notion of edge
weight beyond hop count.

To our knowledge, \Mathlib{} as of 2025 contains no executable, edge-weighted
shortest-path algorithm with a closed correctness theorem. A handful of
community pull requests and \texttt{Zulip} discussions have circled around
adding BFS or DFS over \code{SimpleGraph}, but none has landed at the time
of writing. This is an unsurprising state of affairs: \Mathlib{}'s
combinatorics library has historically been driven by the needs of the
graph-theory and combinatorics formalisation community
(Doczkal--Pous treewidth-two minor exclusion in \textsc{Coq}/MathComp~\cite{DoczkalPous2020Minor}
is a typical comparator), not by algorithm-engineering needs. A natural
home for an algorithmic SSSP development in \Lean{} is therefore a
self-contained namespace --- \code{Sssp.*} in our case --- that re-uses
\Mathlib's structural primitives (\code{Walk}, \code{Path}, the
\code{WithTop}/\code{NNReal} order theory) without yet pushing for upstream
inclusion.

Comparable algorithmic graph developments outside \Lean{} include
Noschinski's Isabelle graph library~\cite{Noschinski2015GraphLib} and the
network-flow line of work in
Isabelle~\cite{LammichSefidgar2016EK,Lammich2019JAR}.

\subsection{Verifying \textsc{Rust}}\label{sec:rel-rust}

The Rust verification ecosystem has matured rapidly in the past five years,
and is now arguably the densest cluster of practical deductive verifiers
for any single language. Six tools merit specific discussion.

\textbf{Verus}~\cite{Lattuada2023Verus,Lattuada2024Verus} embeds
specifications, proofs and executable code in a single Rust dialect via a
mode system (\code{spec}, \code{proof}, \code{exec}); ghost types are
linearly typed, allowing proofs to manipulate ownership of memory and
concurrent resources, and Z3 discharges the verification conditions.
\textbf{Creusot}~\cite{DenisJourdanMarche2022Creusot} translates safe Rust
to \textsc{Why3}/\textsc{Coma} via the \emph{Pearlite} specification
language, encoding mutable borrows by way of prophecy variables;
\textsc{Why3}'s SMT back-ends discharge the bulk of the resulting goals.
\textbf{Aeneas}~\cite{HoProtzenko2022Aeneas,HoFromherzProtzenko2024Aeneas}
takes a different tack: it gives a value-based, ownership-centric pure
$\lambda$-calculus semantics to the LLBC borrow calculus, and emits
functional code in F$^\star$, \textsc{Coq}, \textsc{HOL4}, or \Lean{}~4
that the user can verify with full interactive theorem proving. Termination
of borrows is handled via novel \emph{backward functions}; soundness with
respect to the Rust borrow checker has subsequently been
established~\cite{HoFromherzProtzenko2024Aeneas}.
\textbf{Prusti}~\cite{Astrauskas2019Prusti,Astrauskas2022Prusti} translates
annotated Rust to the \textsc{Viper} intermediate verification language,
exploiting Rust's type-based separation of references for modular reasoning.
\textbf{Kani}~\cite{VanHattum2022Kani} is a bounded model-checker built on
\textsc{CBMC}, targeting panics, integer overflows, and undefined behaviour
rather than full functional correctness.
\textbf{RustHorn}~\cite{Matsushita2020RustHorn} encodes Rust verification
problems as constrained Horn clauses using a prophetic model of mutable
borrows; \textbf{RustHornBelt}~\cite{Matsushita2022RustHornBelt} grounds the
encoding semantically in the \textsc{RustBelt}~\cite{JungRustBelt2018}
framework, extending it to unsafe code.
The \textsc{hax}/\textsc{hacspec} toolchain~\cite{HacspecHax2024} is
crypto-focused but architecturally similar to Aeneas: a Rust subset is
translated to F$^\star$, \textsc{Coq}, \textsc{EasyCrypt}, or
\textsc{ProVerif} for verification.

The striking observation, for the present paper, is that \emph{none} of
these toolchains has, to our knowledge, produced a published verification
of \Dijkstra{} or any non-trivial SSSP algorithm. The closest published
non-trivial verifications are CreuSAT (a verified Rust SAT solver in
\textsc{Creusot}) and the systems-verification line of work on top of
\textsc{Verus} (NrOS, IronKV-style replays, \emph{Atmosphere}). Our work
fills the SSSP gap directly, taking an \textsc{Aeneas}-style
\emph{hand-translated} approach that is methodologically compatible with
future automation through the Aeneas+Charon pipeline.

\subsection{Floating-point in proof assistants and in shortest-path
work}\label{sec:rel-fp}

The state of the art for IEEE-754 floating-point in proof assistants is
\textsc{Flocq}~\cite{BoldoMelquiond2011Flocq,BoldoMelquiond2017Book}: a
\textsc{Coq} formalisation parameterised over the rounding mode and
precision, with a comprehensive theory of representable numbers, ULP
arithmetic, and forward/backward error bounds.
\textsc{CompCert}'s C semantics~\cite{BoldoJourdanLeroyMelquiond2014} is
linked to \textsc{Flocq}, giving an end-to-end pipeline from C source to
machine code that preserves bit-exact FP behaviour. On top of this stack,
\textsc{VST}'s round-off-error tooling
(\textsc{LAProof}/\textsc{VCFloat~2}~\cite{KellisonAppelBertotTassarotti2024VCFloat2})
gives Coq tactics for proving forward error bounds on real C programs,
and Boldo et al.~\cite{BoldoFilliatreMelquiond2009Gappa} have shown how to
interface \textsc{Coq} with the \textsc{Gappa} numerical solver for
near-automatic FP proofs. Ramananandro et
al.~\cite{Ramananandro2016CPP} present an alternative Coq framework for
end-to-end FP-aware C verification.

\Mathlib's IEEE-754 story is, by contrast, very thin.
\code{Mathlib.Data.Float.Basic}~\cite{LeanFloat} is a wrapper around the
runtime \code{double}: it provides the algebraic operations and the
order, but no rounding-mode-parameterised semantics, no
ULP arithmetic, and no formalised IEEE-754 algebra of the kind \textsc{Flocq}
provides. As a consequence, anyone wishing to prove a forward-error bound on
an FP algorithm in \Lean{} must currently either build their own IEEE-754
model from \code{BitVec}, port a fragment of \textsc{Flocq}, or
axiomatise the necessary facts and discharge them out-of-band. Our paper
takes the third option, with the named axiom block
\code{NumericBridge} (\cref{sec:refine}) listing the IEEE-754 facts that
the Dijkstra refinement needs.

The intersection of mechanized FP and shortest-path verification is, to our
knowledge, empty. The numerical-analysis literature does have informal
treatments: Higham's \emph{Accuracy and
Stability}~\cite[Ch.~4]{Higham2002Accuracy} gives the standard
$(n-1)\cdot\mathrm{ulp}$ forward-error bound for an $n$-term summation, which
specialises directly to a shortest-path of $n-1$ edges, but this has not
been mechanized for \Dijkstra{} in any system we are aware of. Our paper
should therefore be read as the necessary first step toward such a
mechanization in \Lean: it puts in place the architectural scaffolding
(\code{Algo}, \code{Refine}, the two remaining bridge axiom modules) on top of which a
\textsc{Flocq}-style forward-error theorem becomes natural future work.

\subsection{Refinement, extraction, and cost analysis}\label{sec:rel-cost}

Stepwise refinement from non-deterministic specification to executable code
has been most thoroughly developed in the Isabelle Refinement
Framework~\cite{NordhoffLammich2012,Lammich2019JAR,Lammich2019LLVM}, where
the AFP \Dijkstra{} entry~\cite{NordhoffLammich2012} is a paradigmatic
example. \textsc{Coq}'s extraction
mechanism~\cite{Letouzey2002Extraction} produces OCaml/Haskell/Scheme code
from \textsc{Coq} terms; \textsc{MetaCoq}~\cite{Sozeau2020MetaCoq} and
\textsc{CertiCoq}~\cite{CertiCoq2017} provide verified-extraction pipelines.
\Lean{}~4 takes a different approach: \code{Decidable} instances yield
executable functions directly, and \code{decide}/\code{native\_decide}
evaluate them at elaboration time, the latter via compiled native code
(with a correspondingly larger trusted base, including the \Lean{} compiler
and runtime). We use \code{native\_decide} only in the regression-fixture
harness (\cref{sec:fixtures}), keeping it strictly outside the trust path of
the main correctness theorem.

Cost analysis in proof assistants has its own substantial line. The canonical
example is Charguéraud and Pottier's union-find with the inverse-Ackermann
amortised bound in \textsc{CFML}/Coq~\cite{CharguraudPottier2015UnionFind,CharguraudPottier2019JAR};
Guéneau, Charguéraud and Pottier~\cite{GueneauChargueraudPottier2018ESOP}
generalise this to a big-$\Bigo$-friendly time-credit framework. Carbonneaux
et al.~\cite{CarbonneauxHoffmannRepsShao2017CAV} produce \textsc{Coq}-checked
cost certificates from \textsc{Why3}, and Hoffmann's \emph{Resource Aware
ML}~\cite{Hoffmann2017RaML} gives an automated analyser with checkable
certificates. \Lean's \code{Asymptotics.IsBigO} (in \code{Mathlib.Analysis})
provides the analytic side of asymptotic notation but has not been linked to
a runtime cost model; building such a cost monad in \Lean{} is left for
future work.

% =====================================================================
\section{Preliminaries}\label{sec:prelim}

This section is a stub; it will fix notation for the rest of the paper.

\paragraph*{Lean 4 and Mathlib.}
We use \Lean{} 4~\cite{deMouraUllrich2021Lean4} with
\Mathlib{}~v4.29.1~\cite{Mathlib2020}. Notation: \code{Fin n} is the finite
type with $n$ elements; \code{WithTop $\alpha$} adjoins a top element
$\top$ to $\alpha$; \code{NNReal} is the type of non-negative reals;
\code{Float} is \Lean's runtime IEEE-754 double-precision wrapper.
\Mathlib's graph theory uses \code{SimpleGraph} but we work with directed
graphs and edge weights, so we introduce our own \code{Graph} record type
in \code{Sssp.Graph}.

\paragraph*{\Dijkstra's algorithm.}
Dijkstra's classical algorithm~\cite{Dijkstra1959} maintains a tentative
distance estimate $\dhat[v]$ for every vertex, initialised to $0$ at the
source $s$ and to $\infty$ elsewhere; it iteratively pops the vertex of
minimum estimate from a priority queue, marks it ``settled'', and relaxes
its outgoing edges $u \to v$ by the rule
$\dhat[v] \leftarrow \min(\dhat[v], \dhat[u] + w(u,v))$. With a binary heap
the running time is $\Bigo((n+m)\log n)$; with a Fibonacci
heap~\cite{FredmanTarjan1987} it is $\Bigo(m + n\log n)$. Standard
textbook treatments give $w \in \mathbb{R}_{\geq 0}$ or
$w \in \mathbb{Z}_{\geq 0}$ as the edge-weight model;
see Cormen et al.~\cite{CLRS2022} or
Sedgewick~\cite{Sedgewick2003GraphAlgorithms}.

\paragraph*{IEEE-754 in two paragraphs.}
Our Refine layer uses \Lean's runtime \code{Float} type, which wraps
IEEE-754 binary64 (\code{f64}) arithmetic~\cite{IEEE754_2019,LeanFloat}.
Arithmetic follows round-to-nearest-even; normal and subnormal values share
the usual exponent--significand layout, and the unit in the last place
(ulp) at exponent~$e$ scales as $2^{e-52}$ for doubles.
We write \code{Float.add}, \code{Float.lt}, and \code{min} for the
operations the Refine code calls; \code{distInf} in
\code{Sssp.Refine.Dijkstra} is \code{1.0 / 0.0}, matching Rust's
\code{f64::INFINITY}.

On the non-negative orthant relevant to nat-cast edge weights
(\code{floatWeight w} $=$ \code{Float.ofNat w}), we rely on monotonicity and
algebraic laws packaged as named axioms in \code{NumericBridge}
(\cref{sec:refine}): addition is associative and commutative on the
operands that arise, and \code{Float.le} is a total preorder compatible
with \code{Float.lt}.
We do \emph{not} prove a forward-error bound in this paper; the standard
Higham-style bound that an $n$-term sum incurs at most
$(n-1)\cdot\mathrm{ulp}$ rounding error~\cite[Ch.~4]{Higham2002Accuracy}
is the natural follow-on once a \textsc{Flocq}-style model is available.
Here we treat IEEE-754 as a first-class specification rather than as a
derived consequence of \code{BitVec} arithmetic; the relevant axioms are
listed and discussed in \cref{sec:refine}.

% =====================================================================
\section{Architecture: spec, algorithm, refinement}\label{sec:architecture}

The \Lean{} development under \code{formal/lean/Sssp/} is organised in three
layers. Each layer answers a different question: what is the correct answer,
what algorithm computes it on abstract graphs, and what executable model
matches the Rust code.

\paragraph*{Three layers.}
\begin{itemize}[leftmargin=1.6em,itemsep=0.25em]
\item \textbf{Spec} (\code{Sssp.<X>})~--- non-computable oracles and metric
      lemmas. Example: \code{trueDist} at
      \code{formal/lean/Sssp/Distance.lean:34}; \code{dijkstraSpec} at
      \code{formal/lean/Sssp/Dijkstra.lean:46}.
\item \textbf{Algo} (\code{Sssp.Algo.<X>})~--- computable algorithms on
      \code{Graph n} with \code{NNReal} weights, proved correct against the
      spec. Example: \code{relaxRound} and \code{Algo.dijkstra} at
      \code{formal/lean/Sssp/Algo/Dijkstra.lean:22--28}.
\item \textbf{Refine} (\code{Sssp.Refine.<X>})~--- operational model with
      \code{Float}, CSR \code{RustGraph}, and lazy heap; mirrors
      \code{src/dijkstra.rs}. Example: \code{dijkstraStep} at
      \code{formal/lean/Sssp/Refine/Dijkstra.lean:326},
      \code{dijkstraHeap} at :366.
\end{itemize}

\begin{verbatim}
  Spec (oracle)          Algo (verified)         Refine (executable)
  ----------------       ----------------        ---------------------
  trueDist,              relaxRound,             floatRelaxRound,
  dijkstraSpec           Algo.dijkstra           dijkstraHeap / dijkstra
         |                      |                         |
         |    dijkstra_correct  |   Simulation + bridges  |
         +----------------------+-------------------------+
                                |
                         refine_dijkstra_correct
                                |
                         src/dijkstra.rs (Rust)
\end{verbatim}

\paragraph*{Why an oracular spec.}
The spec layer factors ``what answer is correct'' from ``what algorithm
computes it''. \code{dijkstraSpec} is defined as \code{trueDist} and
\code{dijkstraSpec\_correct} holds by \code{rfl}; the real work lives in
\code{Sssp.Algo.Dijkstra}. This pattern lets future algorithms (e.g.\ a
verified \BMSSP{}) share the same \code{trueDist} metric without duplicating
distance theory.

\paragraph*{Why CSR and \code{Float} in Refine.}
The Refine layer is hand-translated from \code{src/graph.rs} (CSR layout,
\code{fromEdgeList}) and \code{src/dijkstra.rs} (lazy min-heap, stale-entry
skip, \code{f64} relax). Using \code{Float} and list-based CSR in \Lean{}
makes a refinement theorem against Rust structurally feasible; a future
\textsc{Aeneas} export could replace the hand translation while preserving
the proof architecture.

\paragraph*{Reading map.}
\begin{center}
\small
\begin{tabular}{@{}lll@{}}
\hline
Lean module & Paper section & Rust source \\
\hline
\code{Sssp.Distance} & \cref{sec:prelim}, \cref{sec:algo} & --- \\
\code{Sssp.Dijkstra} & \cref{sec:algo} & --- \\
\code{Sssp.Algo.Dijkstra} & \cref{sec:algo} & --- \\
\code{Sssp.Refine.Dijkstra} & \cref{sec:refine} & \code{src/dijkstra.rs} \\
\code{Sssp.Refine.RelaxBridge} & \cref{sec:refine} & --- \\
\code{Sssp.Refine.GraphBridge} & \cref{sec:refine} & \code{src/graph.rs} \\
\code{Sssp.Refine.NumericBridge} & \cref{sec:refine} & --- \\
\code{Sssp.Refine.HeapBridge} & \cref{sec:refine} & \code{src/dijkstra.rs} \\
\code{Sssp.Refine.RefineCorrectness} & \cref{sec:refine} & --- \\
\code{Sssp.Fixtures.*} & \cref{sec:fixtures} & \code{src/dijkstra.rs} \\
\hline
\end{tabular}
\end{center}

% =====================================================================
\section{The verified abstract \Dijkstra}\label{sec:algo}

The verified algorithm lives in \code{Sssp.Algo.Dijkstra}. It performs
$n$ global relaxation rounds on a directed graph \code{G : Graph n} with
non-negative \code{NNReal} edge weights.

\paragraph*{Definition.}
One round relaxes every vertex's outgoing edges:
\code{relaxAll} folds \code{relaxOutEdges} over \code{List.finRange n};
\code{relaxRound G fuel} recurses on \code{fuel}; and
\begin{verbatim}
def dijkstra (G : Graph n) (s : Fin n) : DistEstimate n :=
  relaxRound G n (initEstimate s)
\end{verbatim}
at \code{formal/lean/Sssp/Algo/Dijkstra.lean:27--28}.
The source estimate sets $\dhat[s]=0$ and $\dhat[v]=\top$ elsewhere
(\code{initEstimate}).

\paragraph*{Soundness.}
The shared lemma \code{relax\_sound}
(\code{formal/lean/Sssp/Dijkstra.lean:57}) shows that if $\dhat$ is
\emph{sound} ($\dhat[v] \geq \code{trueDist}\,G\,s\,v$ for all $v$), then
relaxing any edge preserves soundness. Induction on \code{fuel} yields
\code{relaxRound\_sound}: after any number of rounds, estimates remain
lower bounds on true distance.

\paragraph*{Completeness.}
For a vertex $v$ with finite distance, \code{exists\_shortest\_bounded\_walk}
(\code{formal/lean/Sssp/Distance.lean:264}) supplies a walk of length
$\leq n$ whose weight equals \code{trueDist}. The private lemma
\code{relaxRound\_le\_initWalk} (\code{Algo/Dijkstra.lean:319}) shows that
$n$ rounds of relaxation along such a walk bring $\dhat[v]$ down to the walk
length, hence to \code{trueDist}. Together with soundness,
\code{dijkstra\_le\_trueDist} and \code{dijkstra\_ge\_trueDist} close the gap.

\paragraph*{Closed theorem.}
\begin{verbatim}
theorem dijkstra_correct {G : Graph n} {s v : Fin n} :
    dijkstra G s v = dijkstraSpec G s v
\end{verbatim}
at \code{formal/lean/Sssp/Algo/Dijkstra.lean:479--480}.
Since \code{dijkstraSpec} equals \code{trueDist}
(\code{Dijkstra.lean:46--52}), this is
the statement that $n$-round relaxation computes exact shortest-path
distances on non-negative weighted graphs---the textbook \Dijkstra{}
invariant in a round-based presentation equivalent to the priority-queue
form on these graphs.

% =====================================================================
\section{Refinement to Rust: bridge axioms and discharge routes}\label{sec:refine}

This section connects the verified \code{Algo} layer to the operational
\code{Refine} model that mirrors \code{src/dijkstra.rs}. Two bridge modules
still carry axioms (\code{NumericBridge}, \code{HeapBridge}); two others
(\code{RelaxBridge}, \code{GraphBridge}) are fully proved.

\subsection{The Refine layer}\label{sec:refine-dijkstra}

\code{RustGraph} stores CSR adjacency via \code{head}, \code{edgeTo}, and
\code{edgeW} (\code{Refine/Dijkstra.lean}, matching \code{src/graph.rs}).
Distances are \code{List Float}; \code{floatRelaxEdge} implements
$\dhat[v] \leftarrow \min(\dhat[v], \dhat[u]+w)$.
\code{dijkstraStep} pops the heap minimum, skips stale entries
(\code{distStale}), and relaxes outgoing CSR edges---the same control flow
as Rust.
\code{dijkstraRun} iterates \code{dijkstraStep} for a fuel bound;
\code{dijkstraHeap} is the production entry point;
\code{dijkstraRelax} is the proof-relevant $n$-round float relaxation used
in the simulation pipeline (\code{Refine/Dijkstra.lean:388--394}).

\subsection{The simulation invariant}\label{sec:refine-sim}

Refinement proofs thread a triple invariant \code{SimInv}
(\code{RelaxBridge.lean:21}):
\begin{enumerate}[leftmargin=1.6em,itemsep=0.2em]
\item \code{len}: the float distance vector has length $n$;
\item \code{aligned}: each entry equals \code{nnrealToFloat} of the parallel
      \code{DistEstimate} entry;
\item \code{sound}: the estimate is sound for \code{trueDist} on
      \code{vg.toGraph}.
\end{enumerate}
Per-edge CSR-vs-\code{Graph} alignment is
\code{relaxOutEdges\_eq\_relaxCsrOut} (\code{RelaxBridge.lean:429}).
Round-level lifting is \code{floatRelaxRound\_simInv}
(\code{Simulation.lean:73}), showing $k$ float rounds mirror $k$ abstract
\code{relaxRound} steps while preserving \code{SimInv}.

\subsection{Axiom-free RelaxBridge and GraphBridge}\label{sec:refine-proved}

\textbf{RelaxBridge.} Former axioms on per-edge and per-round alignment are
discharged by length lemmas, \code{floatRelaxEdge\_aligned}, fold induction
on \code{List.range}, and \code{relaxOutEdges\_eq\_relaxCsrOut}. The module
contains no \code{axiom} declarations.

\textbf{GraphBridge.} \code{outEdge\_floatWeight\_preimage} links CSR edge
indices to \code{Graph.outEdges}, using \code{fromEdgeList}'s prefix-sum
construction. Together with fixture \code{ValidRustGraph} proofs, this
connects JSON test graphs to the verified \code{Graph n} type.

\subsection{NumericBridge: 21 IEEE-754 facts}\label{sec:refine-numeric}

\code{NumericBridge.lean:80--122} names 21 axioms in four groups:

\paragraph*{Float order and algebra (9 facts).}
\code{float\_le\_refl}, antisymmetry, transitivity, \code{float\_le\_of\_lt},
\code{float\_le\_of\_not\_lt}, \code{float\_le\_top}; plus
\code{floatZero\_add}, \code{float\_add\_assoc}, \code{float\_add\_comm}.
These would follow from a full IEEE order specification on \code{Float};
\Mathlib{} does not yet supply them for proof.

\paragraph*{\code{floatWeight} $\leftrightarrow$ \code{Float.ofNat} (4 facts).}
\code{floatWeight\_eq\_ofNat}, \code{floatWeight\_add},
\code{floatWeight\_lt\_iff}, \code{floatWeight\_le\_iff} pin nat-cast
weights to exact \code{Float} arithmetic on small integers---the regime our
fixtures and Rust code use.

\paragraph*{\code{nnrealToFloat} bridge (4 facts).}
\code{nnrealToFloat\_monotone}, \code{nnrealToFloat\_trueDist\_add},
\code{nnrealToFloat\_add\_weight}, \code{nnrealToFloat\_min} lift the
\code{NNReal} metric through the float embedding used in alignment.

\paragraph*{\code{min} and \code{BEq} (4 facts).}
\code{float\_min\_eq\_left/right\_of\_le}, \code{float\_eq\_of\_beq} support
target-vertex relaxations where \code{min} selects the improving branch.

\subsection{HeapBridge: one axiom and three conditional routes}\label{sec:refine-heap}

The standing axiom is
\begin{verbatim}
axiom dijkstraHeap_eq_dijkstraRelax (vg : ValidRustGraph n g) (source : Nat) :
    dijkstraHeap g source = dijkstraRelax g source
\end{verbatim}
at \code{HeapBridge.lean:23}. Three theorems discharge it under hypotheses:
\code{dijkstraHeap\_eq\_dijkstraRelax\_of\_upper} (:27),
\code{of\_complete} (:43), \code{of\_edgeUpper} (:55).

The proof route is textbook Dijkstra completeness transported to
\code{dijkstraRun}: with fuel \code{dijkstraHeapFuel g}, show every vertex's
tentative estimate in \code{dijkstraRun\_dHat} is \code{IsComplete}, then apply
\code{dijkstraHeap\_eq\_dijkstraRelax\_of\_complete}.
Supporting lemmas include \code{dijkstraRun\_processed\_card\_le\_fuel},
\code{freshPop\_isComplete\_of\_processed\_pred}, and schedule lemmas
\code{dijkstraRun\_dHat\_schedule\_of\_*} in \code{HeapSimulation.lean}.

\textbf{Open obligation.} Prove
\code{dijkstraRun\_dHat\_all\_complete\_at\_heapFuel}: after heap-loop fuel
expires, every vertex's estimate equals true distance. This is the single
remaining cut for axiom-free \code{HeapBridge}.

\subsection{Headline theorem}\label{sec:refine-headline}

\begin{verbatim}
theorem refine_dijkstra_correct (vg : ValidRustGraph n g) (s v : Fin n) :
    (dijkstra g s.val)[v.val]! = withTopNatToFloat (trueDistNat vg.toGraph s v)
\end{verbatim}
at \code{RefineCorrectness.lean:23--25}. The proof is three steps:
(1)~\code{dijkstra\_get\_eq\_dijkstraRelax} rewrites heap output to the
relaxation model;
(2)~\code{refine\_dijkstraRelax\_correct} connects relaxation to
\code{trueDistNat} via \code{Simulation} and the bridges;
(3)~\code{trueDistNat\_toFloat} links the nat shadow to the float answer.
Step~(2) uses \code{HeapBridge} (axiom) and the axiom-free
\code{RelaxBridge}/\code{GraphBridge}/\code{NumericBridge} chain.

% =====================================================================
\section{Fixture-based regression harness}\label{sec:fixtures}

Operational agreement between Lean and Rust is checked on four shared JSON
graphs under \code{formal/fixtures/dijkstra/}:
\code{tiny\_chain.json}, \code{diamond\_with\_ties.json},
\code{unreachable\_vertices.json}, and \code{single\_vertex.json}.
Each file specifies vertex count $n$, source, weighted edges
\code{(u, v, w)}, and expected distance vector.

On the Lean side, \code{Sssp.Fixtures.Graph} builds \code{RustGraph} values;
\code{Sssp.Fixtures.Dijkstra} runs \code{Refine.dijkstra}; and
\code{Sssp.Fixtures.Correctness} uses \code{\#guard} with
\code{native\_decide} to assert bit-level equality against expected floats.
On the Rust side, \code{cargo test shared\_json\_fixtures} reads the same
paths (\code{src/dijkstra.rs:99--113}), builds \code{Graph::from\_edges}, and
compares \code{dijkstra} output to \code{expected\_dist}.

The glue script \code{formal/scripts/check-fixtures.sh} runs both sides in CI
(see \code{.github/workflows/}). This harness is \emph{regression only}: it
does not appear in the proof of \code{refine\_dijkstra\_correct}.
Using \code{native\_decide} enlarges the trusted base to include the \Lean{}
compiler and runtime for those guards; the symbolic refinement theorem does
not depend on it.

% =====================================================================
\section{Empirical study: \Dijkstra{} versus \BMSSP}\label{sec:benchmark}

We benchmark our Rust \Dijkstra{} (binary heap, \code{f64}, CSR) against an
implementation of \BMSSP{}~\cite{DuanMaoMaoShuYin2025} on Erdős--Rényi-style
random directed graphs: weights uniform in $(0,1]$, source vertex~$0$, average
out-degree~$4$. $L$ is \BMSSP{}'s recursion depth; \code{ratio} is
\BMSSP-time\,/\,\Dijkstra-time; \code{ratio$\cdot$log$^{1/3}$(n)} rescales by
the predicted asymptotic gap.

\begin{center}
\small
\begin{tabular}{@{}rrrrrrrr@{}}
\hline
$n$ & $m$ & $L$ & Dijkstra & \BMSSP & ratio & ratio$\cdot$log$^{1/3}$ & match \\
\hline
1000 & 4000 & 3 & 109.2\,$\mu$s & 554.3\,$\mu$s & 5.08$\times$ & 10.92 & yes \\
3000 & 12000 & 3 & 375.8\,$\mu$s & 1.72\,ms & 4.59$\times$ & 10.37 & yes \\
10000 & 40000 & 3 & 1.59\,ms & 5.90\,ms & 3.72$\times$ & 8.80 & yes \\
30000 & 120000 & 3 & 6.81\,ms & 20.73\,ms & 3.05$\times$ & 7.49 & yes \\
100000 & 400000 & 3 & 34.13\,ms & 85.64\,ms & 2.51$\times$ & 6.40 & yes \\
300000 & 1200000 & 4 & 146.61\,ms & 374.60\,ms & 2.56$\times$ & 6.72 & yes \\
1000000 & 4000000 & 3 & 724.31\,ms & 1.570\,s & 2.17$\times$ & 5.88 & yes \\
1500000 & 6000000 & 3 & 1.113\,s & 2.647\,s & 2.38$\times$ & 6.51 & yes \\
2000000 & 8000000 & 3 & 1.661\,s & 4.010\,s & 2.41$\times$ & 6.65 & yes \\
2500000 & 10000000 & 4 & 2.091\,s & 6.029\,s & 2.88$\times$ & 7.99 & yes \\
3000000 & 12000000 & 4 & 2.508\,s & 7.856\,s & 3.13$\times$ & 8.71 & yes \\
4000000 & 16000000 & 4 & 3.577\,s & 12.559\,s & 3.51$\times$ & 9.83 & yes \\
\hline
\end{tabular}
\end{center}

Three readings follow from the sweep (reproduced from the repository README):
\begin{enumerate}[leftmargin=1.6em,itemsep=0.25em]
\item \textbf{Asymptotic ratio decay.} Within the $L=3$ regime, raw ratio falls
      $5.08\times \to 2.17\times$ while \code{ratio$\cdot$log$^{1/3}$} drops
      $10.92 \to 5.88$, matching the predicted
      $\Theta(m\log^{2/3} n)$ vs.\ $\Theta(m\log n)$ gap.
\item \textbf{Recursion-depth steps.} When $L$ ticks up (e.g.\ near
      $n\approx 300\mathrm{k}$ and past $2\mathrm{M}$),
      \code{ratio$\cdot$log$^{1/3}$} jumps $\sim$25--35\% per extra frame.
\item \textbf{Constants dominate.} At $n=4\cdot 10^6$, \BMSSP{} is still
      $3.5\times$ slower. Crossover needs $\log^{1/3} n \gtrsim 6$, i.e.\
      $n \gtrsim 2^{216}$---far beyond workstation scale.
\end{enumerate}

For this paper's scope, a verified textbook \Dijkstra{} is the practical
baseline; \BMSSP{} appears only as an unverified timing comparison motivating
future work. Every row reports \code{match=yes} (distances agree within
$10^{-7}$ relative tolerance).

% =====================================================================
\section{Discussion and future work}\label{sec:discussion}

\paragraph*{Phase~3c: closing \code{HeapBridge}.}
RelaxBridge and GraphBridge are axiom-free. The remaining heap obligation is
\code{dijkstraRun\_dHat\_all\_complete\_at\_heapFuel}: a textbook Dijkstra
completeness argument on \code{dijkstraRun}, using
\code{freshPop\_isComplete\_of\_processed\_pred} and fuel cardinality lemmas
already in \code{HeapSimulation.lean}. Discharging it yields axiom-free
\code{HeapBridge}; only \code{NumericBridge}'s 21 IEEE-754 facts would remain.

\paragraph*{NumericBridge and \textsc{Flocq}.}
The 21 axioms form a stable, named specification of \code{Float} behaviour on
nat-cast inputs. A meaningful next step is proving the nine order/algebra facts
from \Mathlib's emerging float theory, or porting a \textsc{Flocq}
fragment~\cite{BoldoMelquiond2011Flocq,BoldoMelquiond2017Book} to discharge
groups wholesale.

\paragraph*{Aeneas-driven extraction.}
The hand-translated Refine layer could be replaced by \textsc{Aeneas}
output~\cite{HoProtzenko2022Aeneas,HoFromherzProtzenko2024Aeneas}, gaining a
machine-checked borrow and memory-safety bridge from Rust to \Lean{} while
keeping the spec/Algo proof split.

\paragraph*{Forward-error bounds.}
With a \textsc{Flocq} foundation, a Higham-style $(n-1)\cdot\mathrm{ulp}$
bound~\cite{Higham2002Accuracy} along shortest paths would connect the
IEEE axiom block to quantitative \code{f64} accuracy claims.

\paragraph*{Cost monad and \BMSSP.}
Linking \Mathlib's \code{Asymptotics.IsBigO} to a runtime cost monad (in the
spirit of Charguéraud--Pottier time
credits~\cite{CharguraudPottier2015UnionFind,GueneauChargueraudPottier2018ESOP})
would recover $(n+m)\log n$ for \Dijkstra{} and eventually support verifying
\BMSSP{}~\cite{DuanMaoMaoShuYin2025} in the same three-layer architecture.
The benchmark in \cref{sec:benchmark} shows that, at workstation $n$, verified
\Dijkstra{} remains the algorithm worth deploying.

\paragraph*{Threats to validity.}
Three limitations deserve explicit mention.
(1)~\code{native\_decide} fixture checks trust the \Lean{} compiler/runtime;
the main theorem does not.
(2)~\code{NumericBridge} axioms are unproven IEEE facts, not derived theorems.
(3) The Rust$\leftrightarrow$\Lean{} correspondence is hand-checked, not
extracted; operational agreement on four JSON graphs does not constitute a
full refinement proof of \code{src/dijkstra.rs}.

% =====================================================================
\section*{Acknowledgements}
We thank the Lean community (Zulip) for Mathlib guidance, the Aeneas team for
the extraction roadmap, and the authors of Duan et al.\ for publishing the
\BMSSP{} reference implementation we benchmark against. Funding and canonical
repository URL to be added at camera-ready.

% =====================================================================
\bibliographystyle{alpha}
\bibliography{refs}

\end{document}
